\chapter{Conclusion and Future Work}
\label{ch:conclusion}

\section{Introduction}
\label{sec:conc-intro}

In this chapter, we present our concluding thoughts on the proposed framework. We also address the constraints encountered during our work and highlight future directions for further exploration in this area of research.

\section{Summary of the Study}
\label{sec:summary}

Our thesis introduced a blockchain-based electronic voting system that integrates a multi-protocol cryptographic stack into a unified voting pipeline to provide broad security coverage. The proposed framework addresses key gaps in existing e-voting research by simultaneously achieving ballot privacy, vote validity, voter anonymity, coercion resistance, privacy-preserving aggregation, post-quantum security, and a Caliper-benchmarked Hyperledger Fabric deployment.

Our framework combines Paillier homomorphic encryption~\cite{paillier1999} for encrypted tallying, Pedersen commitments~\cite{pedersen1991} with zero-knowledge $\Sigma$-proofs (Strong Fiat-Shamir~\cite{bernhard2012}) for vote validity verification, linkable ring signatures~\cite{liu2004} for anonymous ballot submission with double-vote detection, SMDC deniable credentials for coercion resistance, SA\textsuperscript{2} samplable anonymous aggregation~\cite{sa2_apple} for two-server vote privacy, Kyber-768~\cite{nist_kyber} hybrid encryption for post-quantum protection, and Damg\aa rd-Jurik-Nielsen threshold decryption~\cite{damgard2010} for distributed key management. To the best of our knowledge, no existing system integrates this many cryptographic protocols in a single, functional voting pipeline.

We implemented the system in Go and deployed it on a Hyperledger Fabric~v2.5~\cite{androulaki2018} permissioned-blockchain network, with the API server's cast endpoint wired through the Fabric Gateway SDK so that each successfully cast vote is published to the deployed chaincode. Performance evaluation confirmed that the proposed framework achieves $O(N)$-per-candidate tally complexity (and $O(N \cdot m)$ total, with $N$ denoting the voter count distinct from the Paillier modulus $n$) with a constant per-vote per-candidate cost of approximately 6.1~$\mu$s, verified empirically from 100 to 100\,000 voters and across $m \in [2, 50]$ candidates. The complete multi-protocol vote casting pipeline reaches a per-vote primitive floor of 68.3~ms (per-phase breakdown sums to 68.27~ms; see Table~\ref{tab:pipeline}), with Pedersen commitments and SMDC credential generation amortised across registration time. Independent Hyperledger Caliper benchmarks of the deployed chaincode in isolation report approximately 200~TPS with a 100\% success rate at 10\,000 transactions; an end-to-end Caliper measurement that drives the integrated cast endpoint with a live Fabric peer in the loop is identified as future work. Post-quantum protection via the Kyber-768 hybrid encryption layer adds only 5.35~ms of overhead per vote (approximately 7.8\% of the 68.3~ms per-vote floor), showing that quantum-resistant security is achievable at modest computational cost.

Our novel SMDC coercion-resistance mechanism, with $k = 5$ credential slots (1 real and 4 fake), provides indistinguishable real and fake slots via Pedersen's perfect-hiding property on the slot commitments. The real-slot index is derived from $\text{HMAC-SHA256}\bigl(k_{\text{srv}},\ \text{``smdc-real-index:''}\,\Vert\,\text{vid}\,\Vert\,\text{eid}\bigr) \bmod k$ using a server-held secret $k_{\text{srv}}$ kept confidential to the election authority; the index is therefore re-derivable only by the authority, is never returned in the credential-issuance response, and is disclosed to the voter only through a duress-bound reveal endpoint that gates the response on a successful match of the voter's registered behavioural signal. This binding denies the index to a coercer present at registration time, and lets a voter under post-registration coercion deceive the coercer with a wrong signal that produces a uniform error response. Under cast-time coercion, voters can additionally present any fake slot together with a wrong behavioural signal to produce a valid-looking but silently discarded vote. SMDC delivers a complete implementation of deniable credentials integrated with a blockchain e-voting pipeline, addressing the implementation gap left open by the theoretical credential-management schemes of JCJ~\cite{jcj2005} and Civitas~\cite{civitas2008} for the works surveyed in Chapter~\ref{ch:related-works}.

Our comparative analysis with state-of-the-art systems confirmed that the proposed framework delivers strong performance while offering broader security coverage. The $O(N)$-per-candidate (and $O(N \cdot m)$-total) tally complexity is a factor-$m$ improvement over ISE-Voting's $O(N \cdot m^{2})$, and the 100\% exact tallying contrasts with BP-Vot's $\geq$98\% differential-privacy approximation. Linear extrapolation of our tally benchmarks projects feasibility for national-scale elections: approximately 5.1~minutes per candidate (and approximately 51~minutes total for $m=10$ candidates) to tally 50~million voters.

\section{Limitations and Future Work}
\label{sec:limitations}

Although our proposed framework has shown notable progress in e-voting security, there are some limitations that can be improved in future work.

\begin{itemize}
    \item \textbf{Universally Composable Security Proofs}: The present thesis provides a two-tier security treatment, Tier~1 game-based reductions of all six theorems to standard hardness assumptions (DCRA, DDH, Module-LWE, DLog) in Chapter~\ref{ch:methodology}, and Tier~2 symbolic-model verification of three protocol-level properties via ProVerif in Chapter~\ref{ch:results}. Section~\ref{subsubsec:uc-composition} additionally outlines the ideal-functionality structure ($\mathcal{F}_{\mathrm{Ballot}}$, $\mathcal{F}_{\mathrm{SMDC}}$, $\mathcal{F}_{\mathrm{SA2}}$) that a future UC formalization should target, but a full Universally Composable (UC) proof~\cite{canetti2001uc,ecclesia2025} is left as future work because it is itself a research-scale undertaking (typically 15--50 pages per protocol primitive in dedicated papers). Future work should formalize the twenty-step pipeline in the UC framework, building on E-cclesia's UC-secure self-tallying construction~\cite{ecclesia2025}, construct UC simulators for each ideal functionality, and apply Canetti's UC composition theorem to combine the per-protocol guarantees into a single concurrent-composition statement. A machine-checked variant in EasyCrypt or CryptoVerif would further strengthen the assurance.
    
    \item \textbf{Full Post-Quantum Migration Roadmap}: Kyber-768 (round-3 CRYSTALS-Kyber via CIRCL~\cite{nist_kyber}) currently protects only the integrity binding of vote ciphertexts; migration to the standardised FIPS-203 ML-KEM-768 implementation is identified as a near-term engineering follow-up. A complete post-quantum migration entails three further replacements. (i) The linkable ring signatures rely on the discrete logarithm problem and should be replaced with a lattice-based linkable ring signature (e.g., one constructed from Module-SIS, the assumption underlying ML-DSA in FIPS~204~\cite{nist_mldsa}). (ii) Paillier homomorphic tallying is vulnerable to Shor's algorithm and should be migrated to a fully homomorphic encryption scheme such as BFV, BGV, or CKKS. (iii) Threshold decryption should likewise be reformulated over the chosen lattice-based scheme. A staged migration is realistic: signature replacement first (ML-DSA is already standardised), homomorphic-tally replacement second (lattice FHE has acceptable performance for the small plaintext domain $\{0, 1\}$ used here), and threshold-decryption replacement last. Until that migration is complete, the present hybrid construction provides confidentiality against harvest-now-decrypt-later adversaries, which is the most pressing of the post-quantum threats for ballots that must remain secret for decades.
    
    \item \textbf{Single Orderer Limitation}: Our current deployment uses a single Raft orderer node, which is sufficient for the development and benchmarking workload reported in this thesis but provides no fault tolerance: any orderer crash halts ordering. Production deployment requires at least three orderer nodes to tolerate one crash failure under Raft, and for high-stakes elections where orderer nodes may also be adversarial, migrating to a Byzantine-fault-tolerant consensus mechanism (e.g., SmartBFT in Fabric~v3.x) would strengthen the system's resilience to malicious nodes.
    
    \item \textbf{Fixed Ring Size}: The ring size is fixed at $|\mathcal{R}| = 100$ in our implementation. While this provides strong anonymity, the signer is statistically indistinguishable among the 100 ring members under DDH, a dynamic ring size mechanism that adapts to the voter pool size and latency requirements would provide more flexible deployment options.
    
    \item \textbf{Blockchain Scalability at Extreme Load}: Our Hyperledger Caliper benchmarks revealed that the single-Raft, two-peer reference configuration shows commit-time degradation at transaction volumes above 25{,}000, attributable to a combination of block-batching at the orderer, version-vector validation on shared election-state reads, and Fabric Gateway connection-pool capacity. Production deployments for national-scale elections would require careful infrastructure engineering including multi-orderer consensus, Gateway connection-pool tuning, horizontal peer scaling, and load balancing.
    
    \item \textbf{Mobile Client Implementation}: Our current implementation provides a server-side API but lacks a dedicated mobile client application. For real-world deployment, a mobile application with on-device biometric authentication, secure enclave key storage, and offline credential management would improve voter accessibility.
    
    \item \textbf{Threshold Parameter Optimization}: The current threshold Paillier key generation takes approximately 30 seconds due to safe prime generation for the DJN scheme. For elections with frequent key rotation, exploring alternative threshold schemes (e.g., based on class groups) could reduce setup time.
    
    \item \textbf{Large-Scale User Study}: Our system has been evaluated through automated benchmarks and unit tests but has not undergone a large-scale user study with real voters. A deployment trial with hundreds of participants would provide valuable data on usability, biometric enrollment success rates, and voter confidence.
\end{itemize}

The limitations we identify above and the future work directions suggest a roadmap for advancing our framework toward production-ready deployment. By addressing formal verification, full post-quantum migration, infrastructure scaling, and user experience design, the system can evolve into a practical solution for secure, coercion-resistant democratic elections.

In summary, the proposed framework achieves the following key quantitative results:

\begin{enumerate}
    \item \textbf{6.1~$\mu$s} per-vote per-candidate tally cost, constant across 100 to 100\,000 voters and across $m \in [2,50]$ candidates ($O(N)$ per candidate, $O(N \cdot m)$ total).
    \item \textbf{68.3~ms} per-vote primitive floor of the multi-protocol cast pipeline (sum of per-phase Table~\ref{tab:pipeline}; Pedersen/SMDC issuance is amortised at registration time and excluded).
    \item \textbf{$\sim$200~TPS} sustained chaincode throughput with 100\% success at 10\,000 transactions, measured by Hyperledger Caliper at the chaincode level (an end-to-end Caliper run through the integrated cast endpoint is left as future work).
    \item \textbf{5.35~ms} Kyber-768 hybrid-encryption overhead per vote ($\approx$7.8\% of the per-vote primitive floor) for post-quantum protection.
    \item \textbf{Multi-protocol cryptographic stack} integrated in a single pipeline, with a broader composition of cryptographic primitives than any system surveyed in Chapter~\ref{ch:related-works}.
    \item \textbf{188 tests} (159 unit/integration + 29 benchmarks) with 100\% pass rate and 78--92\% coverage on core packages.
\end{enumerate}

In summary, we believe our study presents a promising and rigorous approach to designing next-generation blockchain-based e-voting systems. Among the works surveyed in Chapter~\ref{ch:related-works}, no prior framework simultaneously provides ballot privacy, vote validity, voter anonymity, coercion resistance, privacy-preserving aggregation, a post-quantum encryption layer, and a Caliper-benchmarked Hyperledger Fabric deployment with thorough empirical validation. While the encryption layer achieves post-quantum protection via Kyber-768 (CRYSTALS-Kyber round-3 via CIRCL; migration to the FIPS-203 ML-KEM-768 standard is identified as a near-term engineering follow-up), the ring signature and homomorphic tallying components remain classically secure, with lattice-based replacements identified as longer-term future work. We believe that the design principles, cryptographic constructions, and engineering practices presented in this thesis can serve as a foundation for future research in this critical area of democratic technology.
